% !TeX program = xelatex
% Standalone Beamer source. Figures, references and optional speaker notes are embedded.
% Compile with XeLaTeX or latexmk -xelatex. No data access or model training is required.
\documentclass[aspectratio=169,11pt]{beamer}
\usepackage{kotex}
\ifdefined\setsanshangulfont
 \IfFontExistsTF{UnDotum.ttf}{
  \setmainhangulfont{UnDotum.ttf}[BoldFont=UnDotumBold.ttf,
   ItalicFont=UnDotum.ttf,BoldItalicFont=UnDotumBold.ttf,Script=Hangul,Language=Korean]
  \setsanshangulfont{UnDotum.ttf}[BoldFont=UnDotumBold.ttf,
   ItalicFont=UnDotum.ttf,BoldItalicFont=UnDotumBold.ttf,Script=Hangul,Language=Korean]
 }{}
\fi
\usepackage{amsmath,amssymb,mathtools,bm}
\usepackage{booktabs,tabularx,array}
\usepackage{tikz}
\usetikzlibrary{arrows.meta,positioning,calc}
\usepackage{listings,xurl}
\usefonttheme{professionalfonts}
\definecolor{ink}{HTML}{17324D}
\definecolor{accent}{HTML}{0B7A75}
\definecolor{bodyink}{HTML}{253443}
\definecolor{pale}{HTML}{EDF4FA}
\definecolor{rulegray}{HTML}{D6DFE8}
\definecolor{warm}{HTML}{A34D20}
\setbeamercolor{normal text}{fg=bodyink,bg=white}
\setbeamercolor{structure}{fg=ink}
\setbeamercolor{frametitle}{fg=ink,bg=white}
\setbeamercolor{title}{fg=ink}
\setbeamercolor{subtitle}{fg=bodyink}
\setbeamercolor{alerted text}{fg=accent}
\setbeamercolor{footline}{fg=ink}
\setbeamercolor{itemize item}{fg=accent}
\setbeamercolor{enumerate item}{fg=accent}
\setbeamerfont{title}{size=\fontsize{24}{29}\selectfont,series=\bfseries}
\setbeamerfont{subtitle}{size=\large}
\setbeamerfont{frametitle}{size=\fontsize{17}{21}\selectfont,series=\bfseries}
\setbeamerfont{footline}{size=\fontsize{6.8}{8}\selectfont}
\setbeamersize{text margin left=9mm,text margin right=9mm}
\setbeamertemplate{navigation symbols}{}
\setbeamertemplate{itemize item}{\raise0.15ex\hbox{\scriptsize\(\bullet\)}}
\setbeamertemplate{itemize subitem}{\raise0.15ex\hbox{\scriptsize\(\circ\)}}
\setbeamertemplate{frametitle}{
 \vspace{2mm}\insertframetitle\par\vspace{1mm}
 {\color{rulegray}\hrule height0.5pt}\vspace{1mm}}
\setbeamertemplate{footline}{
 \leavevmode\hbox{\begin{beamercolorbox}[wd=\paperwidth,ht=3.4mm,dp=1.4mm,
  leftskip=9mm,rightskip=9mm]{footline}
  \usebeamerfont{footline}PADS PD/ET\quad MultiROCKET--HYDRA
  \hfill\insertframenumber\,/\,\inserttotalframenumber
 \end{beamercolorbox}}}
\setbeameroption{hide notes}
\hypersetup{unicode=true,pdftitle={PADS PD/ET MultiROCKET-HYDRA 한국어 Beamer},
 pdfauthor={PD/ET Classification Project},colorlinks=true,urlcolor=accent,linkcolor=ink}
\pdfstringdefDisableCommands{\def\translate#1{#1}}
\newcommand{\R}{\mathbb R}
\newcommand{\ind}{\mathbf 1}
\newcommand{\PD}{\mathrm{PD}}
\newcommand{\ET}{\mathrm{ET}}
\newcommand{\code}[1]{\texttt{\detokenize{#1}}}
\newcommand{\src}[1]{\par\vfill{\scriptsize\color{gray}#1}\par}
\newcommand{\lead}[1]{{\color{ink}\bfseries #1}\par\vspace{1.5mm}}
\DeclareMathOperator{\BA}{BA}
\DeclareMathOperator{\MR}{MR}
\DeclareMathOperator{\HY}{HY}
\DeclareMathOperator{\diag}{diag}
\DeclareMathOperator*{\argmin}{arg\,min}
\DeclareMathOperator*{\argmax}{arg\,max}
\setlength{\emergencystretch}{1.5em}
\setlength{\parskip}{0.3em}
\renewcommand{\arraystretch}{1.18}
\lstset{basicstyle=\ttfamily\fontsize{9}{11}\selectfont,
 columns=fullflexible,keepspaces=true,showstringspaces=false,
 breaklines=true,breakatwhitespace=false,
 frame=none,backgroundcolor=\color{pale},aboveskip=1mm,belowskip=1mm}
\title{PADS 데이터의 PD/ET 분류}
\subtitle{MultiROCKET--HYDRA\\모델 구조와 평가 설계}
\author{PD/ET Classification Project}
\date{2026년 10월 17일\\[1mm]\small 구현 기준: aeon 1.6.0}

\begin{document}

% 01
\begin{frame}[plain]
\vspace{5mm}
\titlepage
\note{이 자료는 모델과 검증 설계를 설명한다. PD/ET 실측 진단 성능에 대한
결과 보고서와 구분한다. 자세한 수학 설명은 함께 제공한 한국어 설명서에 있다.}
\end{frame}

\section{문제와 데이터}
% 02
\begin{frame}[t]{발표 범위}
\lead{참가자별 움직임으로 공개 PD/ET 라벨을 예측하는 방법}
\begin{enumerate}
 \item 데이터의 단위와 입력 행렬
 \item MultiROCKETㆍHYDRA 특징과 가중 릿지
 \item 참가자 단위 중첩 교차검증
 \item 지표 해석과 재현 방법
\end{enumerate}
\vspace{2mm}
\[
 \text{평가 단위: 참가자}\qquad
 \text{출력: ET를 지지하는 실수 점수}
\]
\src{본문은 발표용 요약이며, 정확한 수치 규약과 증명은 부록에 제시한다.}
\note{핵심 구분은 참가자와 기록 조각, 특징 척도 조정과 릿지 벌점,
평가용 외부 예측과 전체 자료에 적합한 최종 모델이다.}
\end{frame}

% 03
\begin{frame}[t]{참가자와 이진 라벨}
\begin{center}
\begin{tabular}{@{}lrr@{}}
\toprule 공개 진단 문자열 & 참가자 & 새 표적 \(y_i\)\\ \midrule
\code{Parkinson's} & 276 & 0\\
\code{Essential Tremor} & 28 & 1\\ \midrule
합계 & 304 & 양성 클래스는 ET\\ \bottomrule
\end{tabular}
\end{center}
\begin{itemize}
 \item 진단 문자열을 정확히 일치시켜 선택한다.
 \item 원래 숫자 라벨 \(2\)는 다른 감별 진단도 포함한다.
 \item ID는 분할과 파일 연결에만 사용한다.
\end{itemize}
\src{\hyperlink{ref-papers}{PADS v1.0.0}\quad \code{data.py: LABELS, load_cohort}}
\note{공개 PD 배정 중 설명에 ET가 등장하는 8명도 공개 라벨을 유지한다.
해당 문구는 감사용 표시이며 재진단이나 예측 특징으로 쓰지 않는다.
출처: https://physionet.org/content/parkinsons-disease-smartwatch/1.0.0/}
\end{frame}

% 04
\begin{frame}[t]{참가자별 입력 행렬}
\[
 X_i\in\R^{132\times976},\qquad
 X\in\R^{304\times132\times976}
\]
\[
 \underbrace{11}_{\text{과제 구간}}\times
 \underbrace{2}_{\text{손목}}\times
 \underbrace{2}_{\text{센서}}\times
 \underbrace{3}_{\text{축}}=132
\]
\begin{itemize}
 \item 채널 순서는 과제 구간, 왼쪽/오른쪽 손목, 가속도/각속도, 축이다.
 \item 서로 다른 과제의 같은 표본 번호는 서로 다른 절대 시각이다.
 \item 한 참가자의 전체 행렬을 하나의 표본으로 유지한다.
\end{itemize}
\src{입력: 공식 \code{preprocessed/movement/ID_ml.bin}\quad little-endian float32}
\note{132는 동시에 측정한 물리 센서 수가 아니다.
보존 구간은 Relaxed1/2, RelaxedTask1/2, StretchHold, HoldWeight,
DrinkGlas, CrossArms, TouchNose, Entrainment1/2이다.}
\end{frame}

% 05
\begin{frame}[t]{공식 전처리와 시간 단위}
\[
 1024\ \text{표본}\;-\;48\ \text{표본}=976\ \text{표본}
\]
\begin{itemize}
 \item 공식 파일에 과제 분할ㆍ선택과 가속도 중력 추세 제거가 이미 들어 있다.
 \item 가속도 단위는 \(\mathrm g\), 각속도 단위는 \(\mathrm{rad}/\mathrm s\)이다.
 \item 현재 로더는 원시 타임스탬프를 복원하거나 재보간하지 않는다.
 \item 명목 100 Hz와 실제 기록의 균일한 시간 간격은 구분한다.
\end{itemize}
\src{\hyperlink{ref-papers}{PADS 사용 지침}\quad 공식 전처리 코드\quad \code{load_cohort}}
\note{앞 48개 표본 제거를 이진 입력에 다시 적용하지 않는다.
976개 표본을 언제나 정확한 9.76초라고 해석하지 않는다.
로더는 해시, 파일 크기, 진단 일치와 선택된 수치의 유한성도 검사한다.}
\end{frame}

\section{모델 구조}
% 06
\begin{frame}[t]{모델의 전체 구조}
\centering
\begin{tikzpicture}[font=\small,>=Latex,
 nodebox/.style={draw=ink,fill=pale,align=center,minimum height=7mm},
 branch/.style={nodebox,text width=47mm,minimum height=11mm},
 middle/.style={nodebox,text width=58mm},
 every path/.style={line width=0.65pt}]
 \node[middle] (input) {\(X_i\in\R^{132\times976}\)};
 \node[branch,below=3mm of input,xshift=-27mm] (mr)
  {MultiROCKET\\일반 특징 표준화};
 \node[branch,below=3mm of input,xshift=27mm] (hy)
  {HYDRA\\제곱근ㆍ영점 보존 정규화};
 \coordinate (mid) at ($(mr.south)!0.5!(hy.south)$);
 \node[middle,below=3mm of mid] (join)
  {\(z_i=[\widetilde h_i;\widetilde r_i]\in\R^{56896}\)};
 \node[middle,below=3mm of join] (ridge)
  {\(s_i=\beta_0+\beta^\top z_i\)};
 \node[middle,below=3mm of ridge] (out)
  {\(s_i>0:\ET\qquad s_i\leq0:\PD\)};
 \draw[->] (input.south)--++(0,-1.5mm)-|(mr.north);
 \draw[->] (input.south)--++(0,-1.5mm)-|(hy.north);
 \draw[->] (mr.south)--++(0,-1.5mm)-|(join.north);
 \draw[->] (hy.south)--++(0,-1.5mm)-|(join.north);
 \draw[->] (join)--(ridge); \draw[->] (ridge)--(out);
\end{tikzpicture}
\src{자료에 의존하는 통계는 학습 집합에서만 추정한다.}
\note{두 특징 분기와 결합 순서는 현재 model.py의 구현을 따른다.
출처: aeon 1.6.0의 MultiRocketHydraClassifier 및 각 변환 소스.
학습 폴드마다 특징 변환과 정규화를 새로 적합한다.}
\end{frame}

% 07
\begin{frame}[t]{적합 대상과 고정 요소}
\small
\begin{tabularx}{\linewidth}{@{}lXl@{}}
\toprule 요소 & 결정 방법 & 라벨 사용\\ \midrule
MR 커널 계열 & 84개 기본 패턴과 설정값 & 없음\\
MR 임계값 & 학습 신호 응답의 분위수 & 없음\\
HYDRA 커널 & 사전에 정한 시드로 생성 & 없음\\
특징의 척도 & 학습 참가자의 특징 통계 & 없음\\
릿지 계수 & 클래스 가중 제곱 손실 최소화 & 사용\\
\(\alpha\) & 내부 검증의 평균 균형 정확도 & 사용\\ \bottomrule
\end{tabularx}
\vspace{2mm}
\normalsize
\alert{라벨을 사용하지 않는 적합도 각 학습 폴드 안에서 수행한다.}
\src{\hyperlink{ref-code}{구현 근거}\quad \code{FeatureMap.fit_transform, tune_alpha}}
\end{frame}

% 08
\begin{frame}[t]{MultiROCKET의 신호 표현}
\[
 (\Delta X_i)_{c,t}=X_{i,c,t+1}-X_{i,c,t}
\]
\[
 v_t=\sum_{c\in\mathcal C}\sum_{\ell=0}^{8}
       a_\ell U_{c,t+(\ell-4)d}
\]
\begin{itemize}
 \item \(U\): 원 신호 또는 1차 차분. 범위 밖의 값은 \(0\).
 \item 길이 9의 기본 패턴은 \(\binom93=84\)개.
 \item 확장 간격 \(d\)와 채널 집합 \(\mathcal C\)를 바꾸어 응답을 얻는다.
\end{itemize}
\src{\hyperlink{ref-papers}{MultiROCKET}\quad \hyperlink{ref-code}{aeon 1.6.0}}
\note{1차 차분은 실제 시간 간격으로 나누지 않는 표본 차분이다.
기본 커널은 세 위치가 2, 나머지 여섯 위치가 -1이다.
현재 normalise=False는 입력 채널별 단위 분산 변환을 적용하지 않는다.}
\end{frame}

% 09
\begin{frame}[t]{MultiROCKET의 네 가지 요약}
\lead{응답 \(v_t\)와 학습 집합에서 얻은 임계값 \(b\)}
\begin{tabularx}{\linewidth}{@{}lX@{}}
\toprule 통계 & 요약하는 정보\\ \midrule
PPV & 임계값을 넘는 표본의 비율\\
MPV & 선택한 응답의 크기\\
MIPV & 선택한 응답의 평균 표본 위치\\
LSPV & 반응 구간의 지속성과 간격\\ \bottomrule
\end{tabularx}
\vspace{2mm}
\[
 \operatorname{PPV}=\frac{|\{t:v_t>b\}|}{m}
\]
\small
MPV의 부호와 LSPV의 경계 처리는 고정한 구현 규약을 따른다.
\src{\hyperlink{pooling-exact}{부록: 정확한 계산 규약}\quad \hyperlink{ref-code}{수치 구현}}
\end{frame}

% 10
\begin{frame}[t]{HYDRA의 경쟁과 누적}
한 그룹의 \(K=8\)개 커널 응답을 \(v_{k,t}\)라 두자.
\[
 k^+(t)=\argmax_kv_{k,t},\qquad k^-(t)=\argmin_kv_{k,t}
\]
\[
 H^+_k=\sum_t v_{k,t}\ind\{k=k^+(t)\},\qquad
 H^-_k=\sum_t\ind\{k=k^-(t)\}
\]
\begin{itemize}
 \item 최대 응답은 응답값 자체를 더한다.
 \item 최소 응답은 선택된 횟수를 센다.
 \item 원 신호와 1차 차분에 모두 적용한다.
\end{itemize}
\src{\hyperlink{ref-papers}{HYDRA}\quad \hyperlink{ref-code}{HydraTransformer}}
\note{동률이면 PyTorch의 첫 인덱스 규약을 따른다.
최소 응답도 반대 부호 패턴에 대한 강한 일치일 수 있다.}
\end{frame}

% 11
\begin{frame}[t]{기본 설정의 특징 차원}
\begin{tabularx}{\linewidth}{@{}lXr@{}}
\toprule 분기 & 계산 & 특징 수\\ \midrule
MultiROCKET & \(2\times4\times84\lfloor6250/84\rfloor\) & 49,728\\
HYDRA & \(2\times7\times32\times8\times2\) & 7,168\\ \midrule
결합 \(z_i\) & \(p=p_{\MR}+p_{\HY}\) & 56,896\\ \bottomrule
\end{tabularx}
\vspace{2mm}
\begin{itemize}
 \item MR의 6250은 특징 예산 설정이며 독립 학습 커널 수가 아니다.
 \item HYDRA의 전체 64개 그룹을 원 신호와 차분에 32개씩 배정한다.
 \item 특징 수가 늘어나도 참가자는 304명이다.
\end{itemize}
\src{현재 aeon 1.6.0, \(T=976\), 기본 설정에 대한 값}
\end{frame}

% 12
\begin{frame}[t]{두 분기의 정규화}
\begin{columns}[T,onlytextwidth]
\begin{column}{0.48\textwidth}
\lead{MultiROCKET}
\[
 \widetilde r_{ij}=\frac{r_{ij}-\mu_j}{\sigma_j}
\]
\small
\(\mu_j,\sigma_j\)는 학습 평균과 척도이다. 상수 특징의 척도는 1을 쓴다.
\end{column}
\begin{column}{0.48\textwidth}
\lead{HYDRA}
\[
 u_{ij}=\sqrt{\max(h_{ij},0)}
\]
\[
 \widetilde h_{ij}=
 \frac{(u_{ij}-\mu_j^{\HY})\ind\{u_{ij}\neq0\}}{s_j+\varepsilon_j}
\]
\small
학습 집합의 평균, 표본 표준편차와 영점 빈도를 사용한다.
\end{column}
\end{columns}
\vspace{2mm}
\normalsize
검증ㆍ테스트 참가자에도 학습 집합에서 정한 통계를 그대로 적용한다.
\src{\hyperlink{sparse-exact}{부록: HYDRA의 척도}\quad 특징 척도 조정과 릿지 벌점은 서로 다른 연산}
\end{frame}

% 13
\begin{frame}[t]{클래스 가중 릿지 목적함수}
\[
 t_i=2y_i-1,\qquad w_c=\frac{n}{2n_c}
\]
\[
 (\widehat\beta_0,\widehat\beta)=
 \argmin_{\beta_0,\beta}
 \left\{\sum_{i\in\mathcal I}w_{y_i}
 (t_i-\beta_0-z_i^\top\beta)^2+\alpha\|\beta\|_2^2\right\}
\]
\begin{itemize}
 \item \(n=|\mathcal I|\), \(n_c\): 학습 집합의 클래스 \(c\) 인원.
 \item 각 클래스의 총 가중치는 \(n_cw_c=n/2\).
 \item \(\alpha>0\)는 계수 크기를 제어하며 절편 \(\beta_0\)에는 벌점을 주지 않는다.
\end{itemize}
\src{\hyperlink{ref-code}{RidgeClassifier}\quad 손실을 \(n\)으로 나누지 않는 구현 규약}
\end{frame}

% 14
\begin{frame}[t]{결정 점수와 예측 규칙}
\[
 s_\star=\widehat\beta_0+\widehat\beta^\top\phi_{\mathcal I}(X_\star)
\]
\[
 \widehat y_\star=
 \begin{cases}
 1\;(\ET),&s_\star>0,\\
 0\;(\PD),&s_\star\leq0.
 \end{cases}
\]
\begin{itemize}
 \item 임계값 0을 평가 전에 고정한다.
 \item 양의 값일수록 ET를 지지하며, 점수를 확률로 해석하지 않는다.
 \item \(s=0.8\)은 ET 확률 80\%를 뜻하지 않는다.
\end{itemize}
\src{\code{FittedModel.decision_function, predict}\quad 확률이 필요하면 별도의 보정 평가를 설계한다.}
\end{frame}

\section{평가 설계}
% 15
\begin{frame}[t]{외부 5폴드와 내부 3폴드}
\[
 \mathcal T_k=\mathcal D\setminus\mathcal E_k,\qquad
 \mathcal D=\{1,\ldots,304\}
\]
\small
\begin{center}
\begin{tabular}{@{}lccccc@{}}
\toprule 외부 폴드 & 분할 1 & 분할 2 & 분할 3 & 분할 4 & 분할 5\\ \midrule
\(k=1\) & \alert{테스트} & 학습 & 학습 & 학습 & 학습\\
\(k=2\) & 학습 & \alert{테스트} & 학습 & 학습 & 학습\\
\(k=3\) & 학습 & 학습 & \alert{테스트} & 학습 & 학습\\
\(k=4\) & 학습 & 학습 & 학습 & \alert{테스트} & 학습\\
\(k=5\) & 학습 & 학습 & 학습 & 학습 & \alert{테스트}\\ \bottomrule
\end{tabular}
\end{center}
\normalsize
각 \(\mathcal T_k\) 안에서만 내부 3폴드로 \(\alpha\)를 선택한다.
\src{참가자 단위 층화 분할. 반복마다 참가자별 외부 예측은 하나.}
\end{frame}

% 16
\begin{frame}[t]{내부 검증의 선택 규칙}
\[
 \mathcal A=\{10^{-3},10^{-2},10^{-1},1,10,10^2,10^3\}
\]
\[
 A_k(\alpha)=\frac13\sum_{j=1}^{3}
           \BA_{\mathcal V_{kj}}(\alpha)
\]
\begin{itemize}
 \item 각 내부 학습 집합에서 특징 변환과 두 정규화기를 다시 적합한다.
 \item 평균 균형 정확도가 가장 높은 \(\alpha\)를 선택한다.
 \item 최댓값과 차이가 \(10^{-12}\) 이하이면 큰 \(\alpha\)를 고른다.
 \item 외부 학습 전체에 재적합하고 마지막에 외부 테스트를 예측한다.
\end{itemize}
\src{\(\mathcal V_{kj}\): 내부 검증 집합\quad \code{tune_alpha, choose_alpha}}
\end{frame}

% 17
\begin{frame}[t]{클래스별 지표와 기준선}
\[
 \BA=\frac12\left(
 \frac{\mathrm{TP}}{\mathrm{TP}+\mathrm{FN}}+
 \frac{\mathrm{TN}}{\mathrm{TN}+\mathrm{FP}}\right)
\]
\begin{center}\small
\begin{tabular}{@{}lcc@{}}
\toprule & 예측 PD & 예측 ET\\ \midrule
실제 PD & TN & FP\\
실제 ET & FN & TP\\ \bottomrule
\end{tabular}
\qquad
\begin{tabular}{@{}lr@{}}
\toprule 항상 PD라고 예측 & 값\\ \midrule
일반 정확도 & \(276/304\approx90.79\%\)\\
균형 정확도 & \(50\%\)\\
ET 재현율 & \(0\%\)\\ \bottomrule
\end{tabular}
\end{center}
\src{기준선은 클래스 구성에서 계산한 값이며 학습 모델의 실측 성능이 아니다.}
\end{frame}

% 18
\begin{frame}[t]{연속 점수에 의한 순위 지표}
\begin{tabularx}{\linewidth}{@{}lX@{}}
\toprule 지표 & 해석\\ \midrule
AUROC & ET 점수가 PD보다 높아지는 정도. 실수 점수로 계산한다.\\
ET의 AP & 임계값에 따른 정밀도와 재현율의 계단형 요약.\\
AP 기준값 & 전원에게 같은 점수를 주면 \(28/304\simeq0.0921\).\\ \bottomrule
\end{tabularx}
\vspace{2mm}
\begin{itemize}
 \item 통합한 외부 점수에는 서로 다른 적합 모델의 척도 차이가 들어간다.
 \item 폴드별 AUROC/AP와 기술적 평균도 함께 확인한다.
\end{itemize}
\src{\hyperlink{ranking-exact}{부록: 순위 지표의 정의}\quad \code{metrics}}
\end{frame}

% 19
\begin{frame}[t]{외부 평가와 최종 모델}
\small
\begin{tabularx}{\linewidth}{@{}lXX@{}}
\toprule & 외부 평가 & 최종 모델\\ \midrule
목적 & 학습 절차의 예측 성능 추정 & 이후 새로운 참가자 예측\\
특징 적합 & 각 외부 학습 집합 & 전체 304명\\
\(\alpha\) 선택 & 외부 학습 안의 3폴드 & 전체 자료 안의 3폴드\\
평가용 예측 & 외부에 보류한 참가자 & 별도의 독립 테스트 필요\\
저장 & 외부 예측과 지표 & \code{final_model.joblib}\\ \bottomrule
\end{tabularx}
\vspace{2mm}
\normalsize
\code{--fit-final}은 전체 자료의 재선택과 재적합을 추가한다.
\src{전체에 적합한 모델의 학습 참가자 예측은 외부 평가에 포함하지 않는다.}
\end{frame}

\section{재현과 해석}
% 20
\begin{frame}[fragile,t]{실행 설정}
\begin{lstlisting}
.venv-sota/bin/python -m pip install \
  -r sota_pd_et/requirements-lock.txt

.venv-sota/bin/python -m sota_pd_et.evaluate \
  --data-root data/pads/1.0.0 \
  --output artifacts/pads_mr_hydra_run2 \
  --outer-folds 5 --inner-folds 3 --repeats 1 \
  --n-jobs 4 --batch-size 32 --seed 20261017 \
  --fit-final
\end{lstlisting}
\small
Python 3.12 환경을 사용하고 출력 경로는 존재하지 않는 새 이름으로 지정한다.
\src{CPU 구현. \code{batch-size}는 한 번에 변환하는 참가자 수.}
\end{frame}

% 21
\begin{frame}[t]{저장 결과의 확인}
\small
\begin{tabularx}{\linewidth}{@{}lX@{}}
\toprule 파일 & 확인할 내용\\ \midrule
\code{summary.json} & 외부 지표, 반복ㆍ폴드별 요약, 기준선\\
\code{oof_predictions.csv} & 참가자별 외부 점수와 예측\\
\code{fold_metrics.csv} & 각 외부 폴드의 지표와 선택한 \(\alpha\)\\
\code{splits.json} & 정확한 외부ㆍ내부 참가자 배정\\
\code{manifest.json} & 실행 상태, 설정, 버전, 자료 해시\\
\code{alpha_search.csv} & 내부 선택용 지표\\
\code{final_model.joblib} & 전체 자료로 재적합한 선택적 최종 모델\\ \bottomrule
\end{tabularx}
\src{외부 예측으로 성능을 읽고 내부 지표로 선택 절차를 확인한다.}
\end{frame}

% 22
\begin{frame}[t]{해석의 범위}
\begin{itemize}
 \item 각 외부 테스트 폴드의 ET는 5명 또는 6명.
       한 명의 예측 변화가 재현율을 크게 바꾼다.
 \item 폴드ㆍ반복의 표준편차는 기술적 산포이다.
       학습 집합이 겹치므로 그대로 신뢰구간이 되지 않는다.
 \item 공식 표현에는 원 연구의 과제 선택과 전처리가 포함된다.
 \item 일반 시계열 벤치마크의 근거, PD/ET 최고 성능,
       외부 기관의 임상 성능을 구분한다.
\end{itemize}
\src{\hyperlink{ref-papers}{PADS 및 Bake off redux}\quad 새 비교는 동일한 분할과 튜닝 예산으로 수행한다.}
\end{frame}

% 23
\begin{frame}[t]{모델과 평가의 핵심}
\[
 X_i\ \xmapsto{\;\phi_{\mathcal I}\;}\ 
 z_i\ \xmapsto{\;\widehat\beta_0+\widehat\beta^\top z_i\;}\ s_i
\]
\begin{itemize}
 \item 두 특징 추출기로 신호의 모양, 반응 위치와 경쟁 패턴을 요약한다.
 \item 학습 집합에서 척도를 정하고 클래스 가중 릿지로 분류한다.
 \item 외부 참가자를 보류한 상태에서 내부 검증으로 \(\alpha\)를 고른다.
 \item 균형 정확도와 ET 재현율을 중심으로 순위 지표를 함께 읽는다.
\end{itemize}
\src{정확한 계산 규약, 유도 과정과 참고문헌은 부록에 제시한다.}
\end{frame}

\appendix
\section{부록 수학}

% A1 / 24
\begin{frame}[t]{부록 A1\quad 주요 기호}
\small
\begin{tabularx}{\linewidth}{@{}lX@{}}
\toprule 기호 & 정의\\ \midrule
\(i,c,t\) & 참가자, 채널, 채널 내부 표본 인덱스\\
\(\mathcal I,\ n\) & 한 번의 적합에 쓰는 학습 집합과 인원 \(n=|\mathcal I|\)\\
\(y_i,t_i\) & ET를 1로 둔 라벨과 \(t_i=2y_i-1\)\\
\(\phi_{\mathcal I}\) & 특징 추출과 분기별 정규화의 합성\\
\(z_i\in\R^p\) & 결합 특징. 기본 \(p=56896\)\\
\(Z\in\R^{n\times p}\) & 학습 특징을 행으로 나열한 행렬\\
\(W=\diag(w_{y_i})\) & 양의 학습 클래스 가중치를 둔 대각행렬\\
\(\beta,\beta_0,\alpha\) & 릿지 계수, 절편, 정규화 강도\\ \bottomrule
\end{tabularx}
\end{frame}

% A2 / 25
\begin{frame}[t,label=pooling-exact]{부록 A2\quad MultiROCKET 풀링의 정확한 정의}
\small
대상 응답을 \(v_0,\ldots,v_{m-1}\),
\(\mathcal P=\{t:v_t>b\}\), \(q=|\mathcal P|\)라 두자.
\[
 \operatorname{PPV}(v,b)=q/m
\]
\[
 \operatorname{MPV}_{\mathrm{impl}}(v,b)=
 \begin{cases}q^{-1}\sum_{t\in\mathcal P}(v_t+b),&q>0,\\0,&q=0,\end{cases}
\]
\[
 \operatorname{MIPV}(v,b)=
 \begin{cases}q^{-1}\sum_{t\in\mathcal P}t,&q>0,\\-1,&q=0.\end{cases}
\]
비교 조건 \(v_t>b\)와 MPV 합산의 \(v_t+b\)를 구분한다.
\src{\hyperlink{ref-code}{aeon 1.6.0 다변량 루틴}\quad LSPV는 다음 슬라이드.}
\end{frame}

% A3 / 26
\begin{frame}[t]{부록 A3\quad LSPV의 경계 규약}
\small
응답 \(v_0,\ldots,v_{m-1}\)와 임계값 \(b\)를 고정한다.
\begin{enumerate}
 \item \(a\gets0,\ L\gets0\)으로 초기화한다.
 \item \(t=0,\ldots,m-1\)의 순서로,
  \(v_t<b\)일 때에만
  \[
   L\gets\max(L,t-a),\qquad a\gets t
  \]
  갱신한다.
 \item 마지막에 \(L\gets\max(L,m-1-a)\)를 반환한다.
\end{enumerate}
\begin{tabularx}{\linewidth}{@{}lX@{}}
\toprule 응답의 예 & 현재 반환값\\ \midrule
모든 \(v_t>b\) & \(L=m-1\)\\
모든 \(v_t<b,\ m\geq2\) & \(L=1\)\\ \bottomrule
\end{tabularx}
\src{엄격한 비교와 끝점 처리를 포함하는 구현 규약. 단순한 연속 양수 개수와 구분한다.}
\end{frame}

% A4 / 27
\begin{frame}[t]{부록 A4\quad HYDRA 커널과 배분}
\small
\[
 a_\ell\overset{\mathrm{iid}}{\sim}\mathcal N(0,1),\quad
 \bar a=\frac19\sum_{\ell=0}^{8}a_\ell,\quad
 \widetilde a_\ell=
 \frac{a_\ell-\bar a}{\sum_{r=0}^{8}|a_r-\bar a|}
\]
\[
 \mathcal D_{\HY}=\{1,2,4,8,16,32,64\},\qquad
 G_{\rm raw}=G_{\rm diff}=32,\quad K=8
\]
\begin{itemize}
 \item 분모가 0이 아니면 계수 합은 0, 절댓값 합은 1이다.
 \item 각 그룹은 132개 채널에서 인덱스 8개를 복원 추출한다.
 \item 채널 혼합은 통계적 연산이며 공통 물리 단위의 계측량은 아니다.
\end{itemize}
\src{\hyperlink{ref-code}{HydraTransformer}\quad 연속 정규분포에서 분모가 0일 확률은 0.}
\end{frame}

% A5 / 28
\begin{frame}[t,label=sparse-exact]{부록 A5\quad HYDRA 희소 특징의 척도}
\small
\[
 u_{ij}=\sqrt{\max(h_{ij},0)},\qquad n=|\mathcal I|\geq2
\]
\[
 \mu_j^{\HY}=\frac1n\sum_{i\in\mathcal I}u_{ij},\quad
 \rho_j=\frac1n\sum_{i\in\mathcal I}\ind\{u_{ij}=0\}
\]
\[
 s_j=\sqrt{\frac1{n-1}\sum_{i\in\mathcal I}(u_{ij}-\mu_j^{\HY})^2},
 \qquad\varepsilon_j=\rho_j^4+10^{-8}
\]
\[
 \widetilde h_{ij}
 =\frac{(u_{ij}-\mu_j^{\HY})\ind\{u_{ij}\neq0\}}{s_j+\varepsilon_j}
\]
영점은 영점으로 남는다. 출력 평균 0과 분산 1을 보장하는 변환은 아니다.
\src{\hyperlink{ref-code}{\code{_SparseScaler}}\quad 모든 통계는 학습 집합에서만 계산한다.}
\end{frame}

% A6 / 29
\begin{frame}[t]{부록 A6\quad 가중 중심화와 정규방정식}
\small
\[
 \bar z_w=\frac{\sum_{i\in\mathcal I}w_{y_i}z_i}{\sum_{i\in\mathcal I}w_{y_i}},
 \qquad
 \bar t_w=\frac{\sum_{i\in\mathcal I}w_{y_i}t_i}{\sum_{i\in\mathcal I}w_{y_i}}
\]
\[
 (Z_c)_{i,:}=(z_i-\bar z_w)^\top,\qquad(t_c)_i=t_i-\bar t_w
\]
\[
 (Z_c^\top WZ_c+\alpha I_p)\widehat\beta=Z_c^\top Wt_c,\qquad
 \widehat\beta_0=\bar t_w-\bar z_w^\top\widehat\beta
\]
\begin{itemize}
 \item \(I_p\)는 \(p\times p\) 단위행렬. 행 \(i\)는 학습 참가자에 대응한다.
 \item 제곱 손실을 계수로 미분한 1차 조건에서 얻는다.
 \item 현재 가중치에서 \(\bar t_w=0\)이어도 절편은 일반적으로 0이 아니다.
\end{itemize}
\end{frame}

% A7 / 30
\begin{frame}[t]{부록 A7\quad 릿지 해의 유일성}
\small
\[
 A=Z_c^\top WZ_c+\alpha I_p,\qquad\alpha>0
\]
임의의 \(v\in\R^p,\ v\neq0\)에 대해
\[
 v^\top Av
 =\|W^{1/2}Z_cv\|_2^2+\alpha\|v\|_2^2>0.
\]
\[
 \widehat\beta=A^{-1}Z_c^\top Wt_c
\]
\begin{itemize}
 \item \(A\)는 양의 정부호이므로 \(p>n\)이어도 계수 해가 유일하다.
 \item 역행렬 식은 실제 구현에서 역행렬을 직접 계산한다는 뜻은 아니다.
 \item 해의 유일성과 새 참가자에서의 예측 성능은 서로 다른 성질이다.
\end{itemize}
\end{frame}

% A8 / 31
\begin{frame}[t,label=ranking-exact]{부록 A8\quad 순위 지표의 정의}
\small
평가 집합의 ET/PD 인원을 \(N_{\ET},N_{\PD}\), 점수를 \(s_i,s_j\)라 두자.
\[
 \operatorname{AUROC}=
 \frac1{N_{\ET}N_{\PD}}
 \sum_{\substack{i:y_i=1\\j:y_j=0}}
 \left[\ind\{s_i>s_j\}+\frac12\ind\{s_i=s_j\}\right].
\]
임계값을 낮추는 순서로 정밀도 \(P_\ell\), 재현율 \(R_\ell\)를 두고 \(R_0=0\)으로 둔다.
\[
 \operatorname{AP}=\sum_\ell(R_\ell-R_{\ell-1})P_\ell.
\]
\begin{itemize}
 \item 동점인 ET/PD 쌍은 AUROC에 \(1/2\)을 더한다.
 \item AP는 계단형 요약이며 사다리꼴 적분과 항상 같지는 않다.
 \item 일정 점수의 AP는 ET 비율이다. 유한 표본의 무작위 순위는 변동한다.
\end{itemize}
\end{frame}

\section{참고문헌}
% A9 / 32
\begin{frame}[t,label=ref-papers]{참고문헌\quad 데이터와 방법}
\small
\begin{description}
 \item[PADS] PADS v1.0.0. PhysioNet, 2024.\\
  \href{https://physionet.org/content/parkinsons-disease-smartwatch/1.0.0/}{Dataset and usage notes}.
  DOI: \href{https://doi.org/10.13026/m0w9-zx22}{10.13026/m0w9-zx22}
 \item[MultiROCKET] C.~W.~Tan et al., 2022.\\
  \href{https://doi.org/10.1007/s10618-022-00844-1}{Multiple pooling operators and transformations}.
 \item[HYDRA] A.~Dempster, D.~F.~Schmidt, G.~I.~Webb, 2023.\\
  \href{https://doi.org/10.1007/s10618-023-00939-3}{Competing convolutional kernels}.
 \item[Bake off redux] M.~Middlehurst, P.~Sch\"afer, A.~Bagnall, 2024.\\
  \href{https://doi.org/10.1007/s10618-024-01022-1}{Time series classification benchmark}.
\end{description}
\end{frame}

% A10 / 33
\begin{frame}[t,label=ref-code]{참고문헌\quad 구현과 평가}
\small
\begin{description}
 \item[aeon 1.6.0] \href{https://github.com/aeon-toolkit/aeon/blob/v1.6.0/aeon/transformations/collection/convolution_based/_multirocket.py}{MultiROCKET transformation source}\\
  \href{https://github.com/aeon-toolkit/aeon/blob/v1.6.0/aeon/transformations/collection/convolution_based/_hydra.py}{HYDRA transformation source}\\
  \href{https://github.com/aeon-toolkit/aeon/blob/v1.6.0/aeon/classification/convolution_based/_hydra.py}{HYDRA sparse scaler source}
 \item[scikit-learn 1.9.1]
  \href{https://scikit-learn.org/stable/modules/generated/sklearn.linear_model.RidgeClassifier.html}{RidgeClassifier API}
 \item[프로젝트] \code{sota_pd_et/model.py}, \code{evaluate.py}, \code{data.py}\\
  자세한 수학 설명: \code{docs/model_explanation_ko.tex}
\end{description}
\vspace{2mm}
문헌의 개념식과 고정 버전의 수치 규약을 구분한다.
\end{frame}

\end{document}
